% ID: MAT-TRI-EQG-001
% Titolo: Valutare affermazioni goniometriche (3)
% Disciplina: Matematica
% Area: Trigonometria
% Argomento: Equazioni goniometriche
% Classe: Quarta secondaria di secondo grado
% Difficolta: 4
% Tipo: vero_falso
% Tag: matematica, trigonometria, equazioni_goniometriche
% Fonte: verifica_goniometria_avanzata_anonimizzata
% Provenienza: materiale fornito dal docente; testo anonimizzato; soluzione redatta per Exergo

\begin{esercizio}
Vero o falso? Giustifica brevemente. 

\begin{enumerate}
\item L'equazione $\sin 2x=\sin x$ ammette esattamente quattro
soluzioni distinte nell'intervallo $[0;2\pi)$. 

\item L'uguaglianza $\dfrac{1-\cos 2x}{\sin x}=2\sin x$
è verificata per ogni $x\in\mathbb{R}$. 

\item L'equazione $\sin x+\cos x=\sqrt{2}$ ammette esattamente
due soluzioni distinte nell'intervallo $[0;2\pi)$. 

\item Il valore $x=\dfrac{\pi}{3}$ è escluso dal dominio della funzione
$g(x)=\dfrac{\sin 2x-\sin x}{\cos x-\frac12}$. 
\end{enumerate}
Salvo diversa indicazione, lavora in $\mathbb R$; angoli in radianti.
\end{esercizio}

\begin{soluzione}
a) Vero: $\sin x(2\cos x-1)=0$ dà $0,\pi,\pi/3,5\pi/3$.
b) Falso: il denominatore impone $x\ne k\pi$; l'identità vale solo nel dominio.
c) Falso: $\sqrt2\sin(x+\pi/4)=\sqrt2$ dà solo $x=\pi/4$ nell'intervallo.
d) Vero: per $x=\pi/3$ il denominatore è nullo, anche se è nullo il numeratore.
\end{soluzione}
